\section{PEFT for DNNs of Other Applications }
\label{sec:peft application}

In Section~\ref{sec:peft taxonomy}, we outlined four categories of PEFT methods along with their improvements. Nonetheless, our discussion did not fully extend to the utilization or adaptation of PEFT techniques beyond traditional architectures (e.g., LLMs) or standard benchmarks (e.g., the GLUE dataset), where the majority of the discussed PEFT methods are applied.
Therefore, in this section, we will highlight and discuss several most representative works that leverage PEFT strategies for various downstream tasks.
We do not aim to cover all PEFT application scenarios in this section.
Our objective is to showcase the significant influence of PEFT within various research domains and demonstrate how to optimize and tailor general-purpose PEFT methods to achieve enhanced performance in specific models or tasks. 

Typically, fine-tuning happens when adapting a pre-trained backbone model to specialized downstream tasks. 
To this end, this section organizes the discussion around various model architectures, which include: LLM, Vision Transformer (ViT), Vision-Language Alignment Model (VLA), and Diffusion model. Within each architectural category, the discussion is further classified based on different downstream tasks.

\subsection{PEFT for LLMs -- Beyond the Basics}
\label{peft_on_llms}

Instead of common tasks in NLP such as NLU and NLG, PEFT techniques boast a wide array of applications across diverse scenarios. PEFT has been successfully implemented in commonsense question answering~\cite{huang2023mvp, zhao2023knowledgeable}, multi-level implicit discourse relation recognition~\cite{zhao2023infusing}, out-of-distribution detection~\cite{ouyang2023prefix}, privacy protection~\cite{ozdayi2023controlling,xiao2023offsite}, federated learning~\cite{che2023federated}, and social biases mitigation~\cite{li2023prompt}. In this section, we pay more focus on three representative downstream tasks: visual instruction following, continual learning, and context window extension.

\subsubsection{Visual Instruct Following}
Several studies, including VL-BART~\cite{cho2021unifying}, MiniGPT-4~\cite{zhu2023minigpt}, and LLaVA~\cite{liu2023visual}, have successfully extended the capabilities of LLMs, initially designed for pure text, to comprehend and generate responses to visual inputs. These enhanced models, namely visual instruct-following LLMs, can process both images and text to produce textual responses, which can be benchmarked on tasks such as image captioning~\cite{rennie2017self,you2016image,vinyals2016show,hossain2019comprehensive} and visual question answering (VQA)~\cite{wang2017fvqa,wu2017visual,antol2015vqa}.
However, these methods fine-tune the entire LLM to learn the visual representations, which can be inefficient in both time and memory. Therefore, it is natural to apply PEFT techniques in the fine-tuning of visual instruct-following LLMs. 
An earlier work \textbf{VL-Adapter}~\cite{sung2022vl} directly applies several PEFT methods (Adapter~\cite{9-peft-transfer-learning}, Hyperformer~\cite{mahabadi2021parameter} and Compacter~\cite{2-Compacter}) on VL-BART~\cite{cho2021unifying} then benchmarks them on several image-text and video-text tasks. Results show that vanilla adapters are the best among them, which can achieve performance on par with full fine-tuning.
However, considering the functionality gap between the encoders and decoders in VL-BART, directly assigning identical modular modifications will lead to suboptimal performance. Therefore, \textbf{VL-PET}~\cite{hu2023vl} selectively integrates PEFT modules into different components of the encoder and decoder. They also introduce a granularity-controlled mechanism for finer-grained control.

To adapt the recently prevalent LLaMA model, \textbf{LLaMA-Adapter}~\cite{zhang2023llama} prepends a set of learnable prompts (similar to prefix tuning) to the input tokens in LLaMA’s higher transformer layers. To avoid the unstable fine-tuning with large loss values at early training stages, instead of the randomly initialized weights of other PEFT methods, LLaMA-Adapter adopts a zero-initialized attention mechanism, which learns a zero-initialized gating factor to adaptively control the contribution of adaptation prompts to the word tokens. This can maintain the fine-tuning starting point the same as the original model and progressively inject new knowledge into the model, where a similar idea can be found in MEFT~\cite{liao2023make} and LoftQ~\cite{loftq} discussed earlier. To represent visual information, LLaMA-Adapter extracts multi-scale global image features using a CLIP image encoder and then projects them to linguistic embedding space. After that, the feature is element-wisely added onto the adaptation prompts at all inserted transformer layers. LLaMA-Adapter only introduces 1.2M learnable parameters in LLaMA-7B and costs less than one hour for fine-tuning on 8 A100 GPUs.
A following work \textbf{LLaMA-Adapter V2}~\cite{gao2023llama} demonstrates that the simple multimodal fusion in LLaMA-Adapter cannot generalize to more challenging open-ended multimodal reasoning tasks, where the visual cues tend to dominate the adaptation prompts than the language instruction data. To address this, LLaMA-Adapter V2 decouples the learning of instruction-following ability (to generate long language responses) and vision-language alignment to avoid interference between visual and language fine-tuning. Specifically, LLaMA-Adapter V2 sets disjoint parameter groups which are respectively learned from image-text pairs and language instruction data. The visual adaptation prompts are inserted in the early stage of LLM, while the language adaptation prompts remain at the higher transformer layers similar to the LLaMA-Adapter. Additionally, LLaMA-Adapter V2 introduces more learnable parameters and several expert systems (e.g., captioning, detection, and OCR) to enhance multimodal performance. 
\textbf{LayerNorm Tuning}~\cite{zhao2023tuning} adjust only the weights of the LayerNorm within each attention block. This straightforward technique can achieve comparable or even better performance than the
finetuning, while offering about 10× more parameter efficiency than LoRA.


\subsubsection{Continual Learning}

Continual Learning (CL) aims to learn a sequence of new tasks over time within one single model, which has broad application in scenarios such as dialogue systems~\cite{lee2017toward}, information extraction systems~\cite{chang2006survey}, and question answering systems~\cite{yang2019end}.
The main challenge in CL is catastrophic forgetting~\cite{kirkpatrick2017overcoming}. A popular practice, called architecture-based methods, tackles the CL by maintaining task-specific parameters in the model for each new task. Therefore, it's natural to leverage PEFT methods for CL tasks~\cite{madotto2020continual,zhu2022continual,dai2022lifelong,liang2023prompts}. For example, \textbf{AdapterCL}~\cite{madotto2020continual} parameterizes each new task using residual adapters. During testing, since the task-id is not provided, AdapterCL uses an entropy-based classifier to select which adapter to use for accomplishing a specific task. \textbf{CPT} (Continual Prompt Tuning)~\cite{zhu2022continual} trains a soft prompt for each task. Instead of training soft prompts from scratch, CPT proposes a series of techniques (continual prompt initialization, query fusion, memory replay, and a memory-guided technique) to achieve knowledge transfer from preceding and subsequent tasks.
\textbf{O-LoRA} (orthogonal low-rank adaptation)~\cite{wang2023orthogonal} employs a strategy of learning distinct tasks within separate low-rank vector subspaces that are kept orthogonal to each other in order to minimize interference. This approach can effectively reduce catastrophic forgetting during the acquisition of new tasks.




\subsubsection{Context Window Extension}

LLMs are typically trained with a pre-defined context size. For example, LLaMA and LLaMA2 have pre-defined context sizes of 2048 and 4096 tokens, respectively. 
The positional encoding RoPE has weak extrapolation properties~\cite{chen2023extending}, which means the performance drops obviously given an input length exceeds the pre-defined context length.
To solve this, a naive solution is to fine-tune a pre-trained LLM to a longer context. However, this escalates computational costs quadratically with context size, straining memory and processing resources.
To address this, \textbf{LongLoRA}~\cite{longlora} proposes to fine-tune a pre-trained LLM using LoRA to enlarge the context size. To reduce the perplexity gap between LoRA tuning and full fine-tuning, LongLoRA also opens embedding and normalization layers for training. In order to further improve training efficiency in a long context scenario, LongLoRA further introduces a novel shifted sparse attention ($\text{S}^2$-Attn) as an efficient substitute for standard self-attention during training.
A subsequent study \textbf{LongQLoRA}~\cite{yang2023longqlora} combines the advantages of LongLoRA with QLoRA and Position Interpolation~\cite{rotationEmbedding} to save GPU memory. This work successfully extends the context length of LLaMA2-13B from 4096 to 8192 on a single V100 with 32GB memory.
\textbf{LLoCO}~\cite{tan2024lloco} introduces a pipeline that learns contexts offline through the combination of context compression and LoRA. The process begins by compressing documents into compact contexts, then fine-tuning LLM using LoRA on the compacted context to improve the LLM’s ability to accurately extract and utilize information from these compressed representations. During model serving, a standard RAG retriever selects both the compressed document and the most relevant LoRA module, and applies them to the LLM for inference. This approach effectively extends the context window of a 4k token LLaMA2-7B model to handle up to 128k tokens.


In addition to limited training-stage sequence length, real-world system memory constraints introduce another critical bottleneck to the context window. Specifically, the capacity of the KV-cache is curtailed by available system memory. For example, a 30B parameter LLM operating with an input length of 1024 and a batch size of 128 might necessitate up to 180GB for the KV-cache~\cite{zhang2024h2o}, thereby restricting the feasible size of the context window. In response to this, some strategies have resorted to quantizing the KV cache~\cite{flexgen, dettmers2022gpt3}, but quantization will certainly compromise performance. To effectively counteract this issue without significant loss, \textbf{GEAR}~\cite{kang2024gear} presents a novel approach by employing a low-rank matrix to capture the majority of coherent bases of quantization error, complemented by a sparse matrix that addresses errors from outlier entries, thus efficiently minimizing approximation errors.



\subsection{PEFT for ViTs}
\label{peft_on_vit}

ViT~\cite{dosovitskiy2010image} has emerged as a powerful backbone model in the recent computer vision community.
In the ViT model, images are treated as sequences of fixed-size patches analogous to how LLM uses discrete tokens. These patches undergo linear embedding and then receive positional encodings. Subsequently, they are processed through standard Transformer encoders. 
The training of ViT can be supervised~\cite{dosovitskiy2010image,steiner2021train} or self-supervised~\cite{chenempirical,he2022masked}, and ViT can achieve superior performance when training with more data and using larger model size~\cite{dehghani2023scaling}. However, such scaling up inevitably escalates training and storage costs. Therefore, similar to LLMs, PEFT is widely implemented in various downstream tasks, such as dense prediction~\cite{chen2022vision}, continual learning~\cite{wang2022learning, gao2023unified}, deep metric learning~\cite{ren2024learning}. Here, we focus on two typical tasks to showcase the involvement of PEFT: image classification and video recognition.

\subsubsection{Image Classification}

Image classification on targeted visual datasets is a very common demand and has extensive applications, while pre-train then fine-tuning paradigm serves as a widespread strategy. A variety of methods leverage PEFT techniques to achieve efficient model tuning~\cite{jia2022visual,chen2022vision,chen2022adaptformer,jie2022convolutional}.
For instance, \textbf{AdaptFormer}~\cite{chen2022adaptformer} inserts adapter modules in parallel to the FFN of the original ViT model for visual recognition tasks. 
\textbf{VPT} (Visual Prompt Tuning)~\cite{jia2022visual} prepends a small amount of task-specific parameters into the input sequence of each Transformer layer. When applying ViT to downstream tasks, only these added parameters and the classification head are set to trainable. 
\textbf{\cite{yoo2023improving}} notices that compared with supervised ViT, VPT often underperforms with self-supervised ViT. Further analysis demonstrates that different pre-trained methods and downstream tasks have varying degrees of dependency on transformer blocks at different locations. To tackle this issue, the research introduces adaptable gates for ViT blocks. These gates dynamically modulate the contribution of prompt tokens to ViT blocks, allowing for a more targeted adaptation of the model to the task at hand.

\subsubsection{Video Recognition}

Several works consider the more challenging adaptation problem that transfers ViT to downstream tasks that have a much larger domain gap. For example, \textbf{ST-Adapter} (Spatio-Temporal Adapter)~\cite{pan2022st} and \textbf{AIM}~\cite{yang2023aim} both insert adapters layers into pre-trained ViT blocks. Their primary goal is to model spatial-temporal information, thereby enabling efficient adaptation of ViTs from image models to video tasks. Notably, both methodologies have exhibited performance that surpasses traditional full-model fine-tuning approaches.


\subsection{PEFT for VLAs}
\label{peft_on_vla}

Vision-language alignment models (VLA), such as CLIP~\cite{radford2021learning}, ALIGN~\cite{jia2021scaling}, DeCLIP~\cite{li2021supervision}, and FLAVA~\cite{singh2022flava}, are designed to learn a good image and text features which can be aligned within a unified representation space.
Each VLA typically consists of separate image and text encoders that extract respective features. Contrastive learning is leveraged in these models to effectively align the image and text features. 
Fine-tuning is leveraged to improve the performance of VLA in specific datasets or tasks, but fine-tuning the full model is computationally intensive. For instance, fine-tuning CLIP RN50x64 requires a batch size of 32,768 and 18 days of training on 592 V100 GPUs~\cite{radford2021learning}. Moreover, full fine-tuning on smaller datasets often leads to catastrophic forgetting~\cite{kirkpatrick2017overcoming}. In response to these challenges, and drawing inspiration from the success of PEFT techniques in NLP, a range of PEFT strategies have been proposed and implemented in VLA models, such as semantic segmentation~\cite{xu2023side, yu2023convolutions, xu2023bridging}, point cloud understanding~\cite{zhang2022pointclip, zhu2023pointclip, wang2022p2p, huang2023clip2point}, video understanding~\cite{ju2022prompting, ni2022expanding, lin2022frozen}, visual reasoning~\cite{han2023zero, doveh2023teaching}, temporal action detection~\cite{nag2022zero}, to name a few.
This section will focus on one common task that uses VLAs: open-vocabulary image classification.

\subsubsection{Open-vocabulary Image Classification}

In open-vocabulary image classification, earlier works design class-specific prompts, e.g., \emph{a photo of a [CLASS]}, for each category, and rank images based on their similarity to these textual descriptions. \textbf{CoOp} (Context Optimization)~\cite{zhou2022learning} replaces the handcrafted text prompt with learnable vectors, while keeping the entire VLA fixes during training. \textbf{CoCoOp} (Conditional Context Optimization)~\cite{zhou2022conditional} builds on this by tackling CoOp's limitations in generalizing to unseen classes. It introduces a lightweight neural network that generates an input-specific context token, dynamically adapting the prompt based on each image, thereby enhancing generalizability, but at the cost of increased computational demands due to the instance-aware operation.
\textbf{ProGrad}~\cite{zhu2023prompt} addresses the over-fitting risk in CoOp in a few-shot setting by regularizing the soft prompt updates whose gradient is aligned to the general knowledge
only updates the prompt whose gradient is aligned (or non-conflicting) to the general knowledge offered by the original prompt.
\textbf{MaPLe}~\cite{khattak2023maple} notes that existing methods learn prompts either in
the language or in the vision branch of CLIP, which is not efficient in leveraging the multimodal nature of VLAs. To address this, MaPLe proposes branch-aware hierarchical prompts that simultaneously adapt both language and vision branches, and achieves superior performance.
\textbf{TPT} (test-time prompt tuning)~\cite{shu2022test} studies prompt tuning on the fly without additional training samples. Specifically, during inference, TPT first augments the input image into various views, which are then utilized to tune the learnable prompts. The primary training objective is to ensure the VLA can generate consistent responses when faced with these differing views. A following work \textbf{DiffTPT}~\cite{feng2023diverse} further enhances the data diversity of test samples through diffusion models.

In another direction, several studies explore the usage of adapters in VLA. For example, \textbf{CLIP-Adapter}~\cite{gao2023clip} integrates residual-style adapters after CLIP’s text and visual encoders. Therefore, unlike CoOp and CoCoOp, CLIP-Adapter avoids the gradient backpropagation through CLIP’s encoders, leading to reduced computational requirements in terms of both training memory and time. \textbf{Tip-Adapter}~\cite{zhang2021tip} adopts the same design with CLIP-Adapter. Different from CLIP-Adapter, the weights of the adapter are obtained in a training-free manner from a query-key cache model~\cite{orhan2018simple,grave2017unbounded} constructed from few-shot supervisions in a non-parametric manner. As a result, Tip-Adapter exhibits great efficiency compared to CLIP-Adapter’s SGD training process.



\subsection{PEFT for Diffusion Models}
\label{peft_on_diffusion}

Diffusion models~\cite{ho2020denoising,sohl2015deep} are a class of generative models that learn to generate data by transforming random noise into a structured output by a progressive denoising process.
During training, diffusion models learn to reverse the noise added to training data using a denoising network, while in inference, they start from noise, using a denoising network to iteratively create data that mirrors the same distribution as the training examples. Diffusion models have various applications~\cite{han2023contrastive,yang2023diffusion,croitoru2023diffusion,dhariwal2021diffusion,ruiz2023dreambooth}, while the most notable is stable diffusion~\cite{rombach2022high}, which bridges the gap between text and image with its robust capability to generate coherent and contextually relevant images directly from textual descriptions. Numerous studies leverage PEFT techniques to adapt a pre-trained diffusion model for downstream tasks, including accelerating sampling speed~\cite{luo2023lcm,chai2023speedupnet}, text-to-video adaptation~\cite{wu2023tune,xing2023simda}, text-to-3D adaptation~\cite{zeng2023ipdreamer}, etc.
This section mainly focuses on two scenarios: integrating additional input modalities beyond mere text-based conditioning, and customizing content generation based on pre-trained diffusion model.



\subsubsection{Additional Input Control}

To incorporate additional input modalities (e.g., layout, keypoints) while retaining the extensive knowledge in the pre-trained model, \textbf{GLIGEN} introduces a novel approach, which maintains the original model's weights intact and integrates new, trainable gated Transformer layers~\cite{alayrac2022flamingo} that take in the new grounding input. The resulting model can not only accurately represent the grounding conditions but also produce high-quality images. Remarkably, the model can also generalize well to unseen objects during inference.
\textbf{ControlNet} \cite{zhang2023adding} fine-tunes a trainable copy of the encoding layers from Stable Diffusion while locking its pre-trained parameter weights. The fixed original model and the trainable copy are bridged through zero convolution layers. These layers, starting with zero-initialized weights, are designed to progressively adapt during training, ensuring that harmful noise does not affect the pre-trained features of Stable Diffusion at the beginning of training. This refined model is capable of conditioning on a variety of inputs such as Canny edges, Hough lines, user scribbles, human key points, segmentation maps, shape normals, depths, etc.
\textbf{Concept Sliders}~\cite{gandikota2023concept} introduces a plug-and-play LoRA adaptors to allow precise editing of concepts (e.g., age, smiling) within a diffusion model.
\textbf{T2I-Adapter} \cite{mou2023t2i} introduces a lightweight adapter model designed to align external control signals with the internal knowledge of text-to-image diffusion models. This adapter enables precise manipulation through structural control (e.g., sketch, depth map, semantic segmentation map, and keypose), color control (e.g., hue and color distribution), and integrating various controls by composing multiple adapters. 




\subsubsection{Customized Generation}

The effectiveness of text-to-image diffusion models is limited by the user's ability to articulate the desired target through text descriptions. For instance, it is difficult to describe the precise features of an innovative toy car which is not encountered during large-scale model training. Consequently, the objective of customized generation is to enable the model to grasp new concepts from a minimal set of user-supplied images.
\textbf{Textual Inversion}~\cite{gal2022image} addresses this by finding a new pseudo-word $S_*$ (similar to soft prompt discussed in Section~\ref{softprompt}) that represents new, specific concepts in the textual embedding space of pre-trained text-to-image diffusion models. The pseudo-word $S_*$ is optimized via the original optimization goal in diffusion models given a small image set (typically 3-5 images) depicting the concept, and the pre-trained model is left untouched. During inference, $S_*$ can be treated like any other word and composed with other textual queries (e.g., "a photo of $S_*$ on the beach").
\textbf{Custom Diffusion}~\cite{kumari2023multi} tackles a more challenging setting: compositional fine-tuning of multiple concepts. It fine-tunes only the $W_k, W_v$ mapping from text to latent features in attention layers, which yields superior performance in multi-concept learning scenarios. Additionally, during fine-tuning, Custom Diffusion prevents model forgetting by introducing a small set of real images with captions akin to the target, alongside employing augmentation for faster convergence and improved results.
\textbf{IP-Adapter}~\cite{ye2023ip} identifies limitations in current approaches (e.g., ControlNet and T2I-Adapter) which project condition signals into the cross-attention modules. When handling image conditions aiming at controlling content, these methods are unable to generate images faithful to the prompted image. The issue stems from that merging image features and text features within cross-attention layers loses image-specific information, leading to only coarse-grained controllable generation such as image style rather than image content. To overcome this, IP-Adapter introduces a novel decoupled cross-attention mechanism to distinguish between text and image features. IP-Adapter adds an additional cross-attention layer exclusively for image features in each cross-attention layer, and only the parameters of the new cross-attention layers are trained.
